\section*{Capítulo \ref{ch:g5:brain-in-command} --- O cérebro no comando}

\begin{solution}{exo:g5:brain-in-command:1}
O \omterm{def:g5:brain-in-command:system}{cérebro}: recebe todos os informes, toma todas as decisões, emite
todas as ordens. A \omterm{def:g5:brain-in-command:system}{medula espinhal}: o cabo grosso que desce do
\omterm{def:g5:brain-in-command:system}{cérebro} por dentro da \omterm{def:g3:bones-and-muscles:skeleton}{coluna}. Os \omterm{def:g5:brain-in-command:system}{nervos}: os fios ramificados que
levam as mensagens entre o \omterm{def:g5:brain-in-command:system}{cérebro} ou a medula e todos os órgãos.
\end{solution}

\begin{solution}{exo:g5:brain-in-command:2}
Os \omterm{def:g1:the-five-senses:sense}{órgãos dos sentidos} informam ao \omterm{def:g5:brain-in-command:system}{cérebro}; o \omterm{def:g5:brain-in-command:system}{cérebro} decide e compõe
as ordens; os \omterm{def:g5:brain-in-command:system}{nervos} levam as ordens até os \omterm{def:g3:bones-and-muscles:muscle}{músculos}, que se
contraem. Os \omterm{def:g1:the-five-senses:sense}{órgãos dos sentidos} informam, o \omterm{def:g5:brain-in-command:system}{cérebro} decide, os
\omterm{def:g3:bones-and-muscles:muscle}{músculos} obedecem.
\end{solution}

\begin{solution}{exo:g5:brain-in-command:3}
O \omterm{def:g5:brain-in-command:system}{cérebro}, no \omterm{def:g3:bones-and-muscles:skeleton}{crânio}; a \omterm{def:g5:brain-in-command:system}{medula espinhal}, dentro da \omterm{def:g3:bones-and-muscles:skeleton}{coluna} de
ossinhos empilhados.
\end{solution}

\begin{solution}{exo:g5:brain-in-command:4}
A demora entre um sinal e a resposta do corpo --- cerca de um quinto
de segundo para uma tarefa simples numa pessoa treinada, e um pouco
mais nas crianças.
\end{solution}

\begin{solution}{exo:g5:brain-in-command:5}
A \omterm{def:g1:the-five-senses:sense}{pele} do pé informou (tato e dor), o \omterm{def:g5:brain-in-command:system}{sistema nervoso} decidiu a
retirada, e os \omterm{def:g3:bones-and-muscles:muscle}{músculos} da \omterm{def:g1:my-body:parts}{perna} obedeceram puxando o pé num tranco.
(Uma retirada tão rápida é resolvida em parte pela própria \omterm{def:g5:brain-in-command:system}{medula
espinhal} --- o balcão de emergência do centro de comando.)
\end{solution}

\begin{solution}{exo:g5:brain-in-command:6}
Ele mantém a \omterm{def:g4:breathing:movements}{respiração} funcionando dia e noite e ajusta o ritmo do
\omterm{def:g4:heart-and-blood:heart}{coração} às necessidades do corpo (e também: guarda lembranças,
sonha).
\end{solution}

\begin{solution}{exo:g5:brain-in-command:7}
O \omterm{prop:g1:healthy-habits:needs}{sono} --- o tempo de arquivamento e reparos dele --- e uma fatia
grande e constante da \omterm{prop:g3:balanced-diet:jobs}{energia} e do oxigênio do corpo. Privado de
qualquer um dos dois, ele reage mais devagar, espalha a atenção e
decide mal.
\end{solution}

\begin{solution}{exo:g5:brain-in-command:8}
Resposta aberta. As pegadas típicas ficam por volta de $10$ a $20$
cm; várias tentativas costumam melhorar e estabilizar um pouco esse
valor.
\end{solution}

\begin{solution}{exo:g5:brain-in-command:9}
Porque toda pegada precisa percorrer os três passos --- o informe
dos \omterm{def:g1:the-five-senses:sense}{olhos} precisa viajar, o \omterm{def:g5:brain-in-command:system}{cérebro} precisa calcular e decidir, a
ordem precisa voltar viajando e os \omterm{def:g3:bones-and-muscles:muscle}{músculos} precisam se contrair ---
e cada passo leva tempo de verdade. A habilidade encurta o circuito;
nada elimina o circuito.
\end{solution}

\begin{solution}{exo:g5:brain-in-command:10}
Porque os melhores freios só agem depois que o circuito de quem
dirige percorreu tudo --- e, em velocidade de rodovia, o carro
percorre uns \qty{7}{m} antes de a ordem ``frear!'' sequer chegar ao
pé. O espaço entre os carros é espaço para o circuito. A atenção
encurta os passos de informar e decidir; um \omterm{def:g5:brain-in-command:system}{cérebro} distraído
acrescenta metros.
\end{solution}

\begin{solution}{exo:g5:brain-in-command:11}
A parte do \omterm{def:g5:brain-in-command:system}{cérebro} --- o decidir e o ordenar: a prática traçou
caminhos melhores, então os mesmos informes agora chegam às ordens
certas depressa e sem esforço. Treinar \omterm{def:g3:bones-and-muscles:muscle}{músculo} aumenta a força do
\omterm{def:g3:bones-and-muscles:muscle}{músculo}; treinar habilidade aumenta o roteamento do \omterm{def:g5:brain-in-command:system}{cérebro}. Um
constrói o motor, o outro constrói o comandante.
\end{solution}
